\section{The complementary three quartic reflection-field class certificates}
\label{sec:cubic-three-ramified-quartic-class-certificates}

Put $p=269$, $t=\sqrt[3]{p}\in\mathbb R$, $K=\mathbb Q(\sqrt{-3})$,
and $z=\zeta_5$.  For $d_0\in\{-11,2,-22\}$ set $d=pd_0$ and
\[
 w_d=\sqrt d\,(z-z^{-1}),\qquad F_d=\mathbb Q(w_d),\qquad
 L_d=\mathbb Q(t,w_d),\qquad M_d=KL_d.
 \tag{G137.1}
\]
Write $D_0=-11,8,-88$ respectively, the fundamental discriminant
of $\mathbb Q(\sqrt{d_0})$.  The fundamental discriminant of
$\mathbb Q(\sqrt d)$ is $pD_0$.
Let $\psi$ be the primitive quartic character modulo five with
$\psi(2)=i$, let $\delta=\psi^2$, and put
$\chi_d=\psi\tau_{D_0}\tau_p$.

\begin{theorem}[AC-G137-ramified-quartic-classes: the complementary class sectors]
\label{AC-G137-ramified-quartic-classes}
For $d=-2959,538,-5918$ one has
\[
                    \operatorname{Cl}(L_d)[269]=0.
 \tag{G137.2}
\]
The nontrivial cubic-character sector of
$\operatorname{Cl}(M_d)[269]$ is consequently zero.
No complete class number, complete unit basis, vanishing of the
trivial cubic sector, or subsequent host/model construction is asserted.
\end{theorem}

\begin{proof}
\proofstep{1. The actual fields and the shared-prime discriminant.}
The two fields $\mathbb Q(z)$ and $\mathbb Q(\sqrt d)$ are disjoint.
In their $C_4\times C_2$ compositum the stabilizer of $w_d$ is
$\{1,s^2r\}$, where $s(z)=z^2$ and $r(\sqrt d)=-\sqrt d$.
It is the kernel of $\chi_d$.  Thus $F_d$ is cyclic quartic,
with characters $1,\delta,\chi_d,\bar\chi_d$, and
\[
 w_d^4+5d w_d^2+5d^2=0,\qquad
 w_d^2=-d(5\pm\sqrt5)/2.
 \tag{G137.3}
\]
The character $\chi_d$ is primitive of conductor $5p|D_0|$;
its parity is even for $d<0$ and odd for $d=538$.
The conductor--discriminant formula gives
$D(F_d)=5^3p^2D_0^2$.
In particular it is unramified at $3$ but has ramification index
two at $p$: its local ramified character there is $\tau_p$.

As in the direct calculation of AC-G135, $T=\mathbb Q(t)$ has
discriminant $-3p^2$.  Indeed it is tamely totally ramified of
degree three at $p$; at $3$, Hensel's lemma applied to
$v+3v^2+3v^3=-30$ gives a cube root $-(1+3v)$ of $p$ in
$\mathbb Q_3$, and the other local factor is
$\mathbb Q_3(\sqrt{-3})$.
There is no other ramification by the polynomial discriminant.
The degrees three and four imply $T\cap F_d=\mathbb Q$.

At every rational prime other than $p$, at least one of the
integral local algebras of $T$ and $F_d$ is unramified.
Their tensor algebra is therefore normal, and the discriminant
exponents are the usual unramified-base-change multiples.
At $p$ this coprime-discriminant shortcut is unavailable.
Both local extensions are tame, with ramification indices three
and two.  Their compositum has ramification index six.
For each completion the tame different exponent is $6-1$;
since the total degree is twelve, the sum of residue degrees is two.
Thus the discriminant exponent at $p$ is ten, rather than fourteen.
It follows that
\[
 \begin{split}
 D(L_d)&=3^4 5^9p^{10}|D_0|^6,\\
 (r_1,r_2,r)&=
 \begin{cases}(4,4,7),&d=-2959,-5918,\\(0,6,5),&d=538.
 \end{cases}
 \end{split}
 \tag{G137.4}
\]
The signatures follow directly from (G137.3).
This is a local field-discriminant calculation, not an identification
of a polynomial-order discriminant with the field discriminant.

The splitting field of $X^3-p$ has group $S_3$ and unique quadratic
subfield $K$.  Its intersection with $F_d$ is trivial, since the
only quadratic subfield of $F_d$ is $\mathbb Q(\sqrt5)$.
Thus $\operatorname{Gal}(M_d/\mathbb Q)=S_3\times C_4$, and $L_d$
is fixed by one reflection of $S_3$.
That reflection normally generates the $S_3$ factor; hence every
abelian subextension of $L_d$ lies in $F_d$.
For $d<0$, $F_d$ is real.  For $d=538$, its unique quadratic
subfield excludes cyclotomic roots of degrees two and four except
possibly fifth roots; those would give $F_d=\mathbb Q(\zeta_5)$,
contrary to the ramification at $p$.
Consequently
\[
                         |\mu_{L_d}|=2.
 \tag{G137.5}
\]

\proofstep{2. Nineteen exact units and their 269-prime indices.}
Use the basis $t^a w_d^b$, $0\le a<3$, $0\le b<4$, at array
index $4a+b$.  Each field's retained
\texttt{all\_units\_exact\_certificate.json} contains respectively
seven, five, seven rational arrays, common denominators, and
monic degree-twelve integral annihilating polynomials with
constant coefficient $\pm1$.
Their exact hashes, in the order $d=-2959,538,-5918$, are
\begin{gather*}
\texttt{b4bd850c001fa9c1f9c3a09be658cd1a}\allowbreak
\texttt{0ecf25297d1a00976c3a285d6bdeb1eb},\\
\texttt{6bbb6109ed4ae237474091f1e063d750}\allowbreak
\texttt{708fd7fa8ae71746efeaf7d8b8d6b2ef},\\
\texttt{16b27fffff294d46eb4b2d38079f6237}\allowbreak
\texttt{ef4284f190168fe687def9749f7c51ad}.
\end{gather*}
All arrays and polynomial coefficients are exact integers divided
by the stated denominators, and are part of the certificate.
They are not obtained by rounding approximate algebraic numbers.

The verifier multiplies in the rational algebra using only
$t^3=p$ and (G137.3), evaluates the monic polynomials by Horner's
rule, and obtains zero exactly for all nineteen elements.
The constant coefficient proves their inverses integral as well.
Thus they are global units.
The raw candidate coordinate is $w_d-t$; both inverse maps to
$t,w_d$ and all conversions to the displayed arrays are also
checked exactly.  The candidate class group and its unit-basis
completeness are not used.

As in AC-G135, take the rational prime $\ell=40351$, primitive
root $3$, and $h=3^{150}=10281\pmod\ell$, of order $269$.
The same trial-division and powering checks apply.
There are three simple roots of $t^3=p$ and four simple roots
of (G137.3) modulo $\ell$, giving twelve embeddings into
$\mathbb Q_\ell$ after Hensel lifting.
They induce residue maps of the full integer ring, since
algebraic integers remain integral under these embeddings.
All displayed unit denominators are invertible modulo $\ell$.
Write $\bar u_i^{150}=h^{c_{ji}}$ under the $j$th residue map.
Every entry is checked by exact powers, and the first $r$ rows
of the complete twelve-row matrix have determinants
\[
\begin{array}{c|r|r}
d&\det(c_{ji})&\det\bmod269\\\hline
-2959&6991745957625888&77\\
538&6018195420&183\\
-5918&6117011886560172&240.
\end{array}
\tag{G137.6}
\]
Let $U_d$ be the group generated by these units and $-1$.
The residue-log homomorphism kills torsion and $269$th powers,
and the matrix has full rank $r$ modulo $269$.
The unit theorem gives rank $r$ for the full unit group modulo
torsion.  Thus the displayed units have full rank and their index
satisfies
\[
 m_d=[\mathcal O_{L_d}^{\times}:U_d]<\infty,
                   \qquad 269\nmid m_d.
 \tag{G137.7}
\]
For clarity, a rational relation can be cleared and divided until
its integer coefficients are not all multiples of $269$, which
would contradict (G137.6).  Once the ranks agree, that same matrix
proves the index prime to $269$ by reduction modulo $269$.

\proofstep{3. Weighted regulator normalization.}
Put $\rho=p^{1/3}$, $\omega=e^{2\pi i/3}$, and
$a_d=\sqrt{|d|(5+\sqrt5)/2}$,
$b_d=\sqrt{|d|(5-\sqrt5)/2}$.
For $d<0$, the four pairs $(\rho,\pm a_d),(\rho,\pm b_d)$
are real embeddings, with weight one, and
$(\rho\omega,\pm a_d),(\rho\omega,\pm b_d)$ represent the
four complex pairs, with weight two.
For $d=538$, the six representatives
$(\rho\omega^j,ia_d),(\rho\omega^j,ib_d)$, $j=0,1,2$,
all have weight two.
Form the weighted-log matrix $q_v\log|\sigma_v(u_i)|$, omit
the last row, and take its absolute determinant $R_{U_d}$.
The projected weighted-log convention gives
\[
                         R_{U_d}=m_dR_{L_d}.
 \tag{G137.8}
\]
There is no Euclidean covolume correction.
Evaluation of these exact formulas at $4000,6000,12000$ bits
respectively gives strictly positive determinants with
\[
\begin{array}{c|c}
d&\text{strict interval containing }R_{U_d}\\\hline
-2959&(3040445199320,3040445199321)\\
538&(5129704519,5129704520)\\
-5918&(180862783024846,180862783024847).
\end{array}
\tag{G137.9}
\]
All weighted-log rows are retained, and each column sum encloses
zero.  The candidate numerical regulators are not inputs to this
calculation, and cancellation is included in the ball errors.

\proofstep{4. The exact Artin factors and local conductors.}
Let $\varrho$ be the standard representation of the $S_3$ field.
Its coefficient sequence $a_n$ is the one in (TQC.19): it is
multiplicative, has $a_{3^v}=1$, $a_{p^v}=0$ for $v\ge1$,
and at a prime $\ell\ne3,p$ satisfies
\[
 a_\ell=\#\{x\bmod\ell:x^3=p\}-1,
 \qquad
 a_{\ell^{v+1}}=a_\ell a_{\ell^v}-(-3/\ell)a_{\ell^{v-1}}.
 \tag{G137.10}
\]
The symbol at $2$ is the Kronecker symbol.
These formulas follow from the actual unramified permutation
representation and the two exceptional local algebras in Step 1.
The local coefficients are bounded by $v+1$, so every twist
satisfies $|b_n|\le d(n)\le n$.

The permutation representation for $L_d$ gives, including all
ramified inertia invariants,
\[
 \operatorname*{res}_{s=1}\zeta_{L_d}(s)
 =L(1,\delta)|L(1,\chi_d)|^2
     L(1,\varrho)L(1,\varrho\delta)|L(1,\varrho\chi_d)|^2.
 \tag{G137.11}
\]
Here $\delta$ is primitive of conductor five; neither primes of
$D_0$ nor $p$ are deleted from the quadratic-twist factor.
The sequences to evaluate are $a_n$, $a_n\delta(n)$ and
$a_n\psi(n)\tau_{D_0}(n)\tau_p(n)$.

Write $\varrho=\operatorname{Ind}_K^{\mathbb Q}\nu$ for the
cubic Kummer character.  Its conductor over $K$ is $(p)$.
At $p$, the norm pullback of $\tau_p$ has nontrivial order-two
residue character, since the residue norm
$\mathbb F_{p^2}^{\times}\to\mathbb F_p^{\times}$ is onto.
The cubic character has order three there.  Their product has
order six, not order one, and remains tame of conductor exponent
one over $K_p$.  This is the shared-prime conductor calculation.

At $5$ the norm pullback of $\psi$ is primitive of exponent one,
and at $11$ the pullback of the quadratic character is primitive
of exponent one when present.  At $2$ the exponent is three
when $8\mid D_0$: in the unramified quadratic local ring,
\[
 N(1+4a)\equiv1+4\operatorname{Tr}_{\mathbb F_4/\mathbb F_2}(\bar a)
                      \pmod8,
 \tag{G137.12}
\]
and the trace is onto while $\tau_8(5)=-1$.
The other factors are unramified at that place and cannot cancel it.
At $3$ all twists are unramified.
The actual Hecke moduli and rational conductors are therefore
\[
\begin{array}{c|c|c}
\text{factor}&f\mathcal O_K&Q=3f^2\\\hline
\varrho&p\mathcal O_K&3p^2\\
\varrho\delta&5p\mathcal O_K&75p^2\\
\varrho\chi_d&5p|D_0|\mathcal O_K&75p^2|D_0|^2.
\end{array}
\tag{G137.13}
\]
The last modulus contains only one factor $p$, despite the
ramification of both original factors at that prime.
All nontrivial ramified local characters have no invariant vector;
the displayed zero-extended coefficient sequences therefore retain
the correct primitive Euler factors.

\proofstep{5. Dirichlet values, including the larger finite modulus.}
The conductor of $\chi_d$ is now $q=5p|D_0|$, namely
$14795,10760,118360$ in the three cases.
The exact finite Fourier computation modulo five gives
\[
                L(1,\delta)=2\log((1+\sqrt5)/2)/\sqrt5.
 \tag{G137.14}
\]
The sign is fixed by
$2\cos(2\pi/5)-2\cos(4\pi/5)=\sqrt5$.
For $d=538$ the primitive character is odd, and exact integer
summation gives
\[
 B_{1,\chi_d}=q^{-1}\sum_{a=1}^{q-1}a\chi_d(a)=-20+12i,
 \qquad |L(1,\chi_d)|^2=544\pi^2/10760.
 \tag{G137.15}
\]
The numerator sums are $-215200$ and $129120$.
The primitive odd functional equation and Gauss norm give the
second formula, with no phase choice.
For $d=-2959,-5918$, put
\[
 S_d=\sum_{a=1}^{q-1}\chi_d(a)\log(2\sin(\pi a/q)),
 \qquad |L(1,\chi_d)|^2=|S_d|^2/q.
 \tag{G137.16}
\]
This is the primitive even finite Fourier identity from AC-G135:
multiply the Abel-convergent cosine series of
$-\log|1-re^{i\theta}|$ by $\chi_d$, sum, and let $r\uparrow1$.
One obtains $S_d=-\mathfrak g(\chi_d)L(1,\bar\chi_d)$ and
$|\mathfrak g(\chi_d)|^2=q$.
All sine arguments lie strictly between zero and $\pi$.
The exact character tables and positive-real logarithms are used
in the ball sums.  No approximation to a Gauss phase occurs.

\proofstep{6. Primitive theta series and a fixed, certified phase.}
The general prime-power argument in AC-G135 applies literally to
all three moduli in (G137.13); we record its precise ingredients.
The field $K$ has class number one by the bound
$2\sqrt3/\pi<2$.
The finite Hecke character $\nu_f$ is primitive modulo
$f\mathcal O_K$, trivial on the six global units, and zero at
nonunits.  For the cubic factor this follows from $9\mid p+1$
and its odd order; norm pullbacks kill these units automatically.
Define $a=2\pi/(\sqrt3f)$ and
\[
 \Theta(t)=\frac16\sum_{\alpha\in\mathcal O_K}
       \nu_f(\alpha)e^{-a|\alpha|^2t}
       =\sum_{n\ge1}b_n e^{-ant}.
 \tag{G137.17}
\]
The inverse-different additive pairing is
$\exp(2\pi i\operatorname{Tr}(\alpha\bar\beta/(f\sqrt{-3})))$.
At a unit frequency its sum is $G\bar\nu_f(\bar\beta)$.
At a nonunit frequency at $\mathfrak q^a$, primitivity provides
a unit in $1+\mathfrak q^{a-1}$ with nontrivial character for
$a>1$, or a residue unit for $a=1$.  Taking it to be one at the
other prime factors leaves the additive character fixed and
multiplies the Fourier sum by a nontrivial scalar.  Thus the
sum is zero.  Parseval gives $|G|=f$.
This includes the factor $2^3$ and the order-six factor at $p$.

The lattice $f\mathcal O_K$ has area $f^2\sqrt3/2$ and dual
$2/(f\sqrt{-3})\mathcal O_K$.  The identity
$\operatorname{Re}(\alpha\overline{(2/\sqrt{-3})\beta})
=\operatorname{Tr}(-\alpha\bar\beta/\sqrt{-3})$
keeps the additive pairing primitive at $2$, without inverting
two in the residue ring.
The Gaussian Fourier transform and these constants give
\[
       \Theta(t)=\varepsilon t^{-1}\bar\Theta(1/t),
                   \qquad |\varepsilon|=1.
 \tag{G137.18}
\]
For $\operatorname{Re}s>1$, Mellin integration equals
$\Gamma(s)a^{-s}L(s)$ by absolute convergence.
The zero constant terms and (G137.18) give exponential decay at
both ends, continuing this equality to $s=1$.
Splitting at one and substituting $t\mapsto1/t$ yields
\[
 L(1)=\sum_{n\ge1}\frac{b_ne^{-an}}n
       +\varepsilon a\sum_{n\ge1}\bar b_n E_1(an),
 \qquad E_1(x)=\int_1^\infty e^{-xt}\frac{dt}{t}.
 \tag{G137.19}
\]
Every computed theta ball at one excludes zero, so the actual
phase is $\varepsilon=\Theta(1)/\overline{\Theta(1)}$.
It equals one for the two real sequences and is evaluated as
a ball quotient for the nonreal sequence.

\proofstep{7. Absolute tails, class-number denominator bound, and isolation.}
For $z_a=e^{-a}$ and $H_N=e^{-a(N+1)}$, the respective omitted
tails of the first sum in (G137.19), its second sum before the
factor $a$, and the theta sum are bounded by
\[
 \frac{H_N}{1-z_a},\qquad
 \frac{H_N}{a(1-z_a)},\qquad
 \frac{H_N((N+1)-Nz_a)}{(1-z_a)^2}.
 \tag{G137.20}
\]
These follow from $|b_n|\le n$ and $E_1(x)\le e^{-x}/x$.
Each real and imaginary component receives the full positive
tail radius.  The calculation uses $256$ bits and the cutoffs
$N_0=12000$, $N_5=50000$, and $N_d=50000|D_0|$.
The regulator balls retain their separately higher precision.

For $n=12$, the exact denominator bound of AC-G135 is
\[
 U=\left\lceil\frac{12!4^{r_2}(\lfloor\sqrt{D(L_d)}\rfloor+1)}
                         {12^{12}3^{r_2}}\right\rceil,
 \qquad H=4^{12}\left\lceil\sqrt[3]{U^4}\right\rceil.
 \tag{G137.21}
\]
Indeed Minkowski supplies a representative of norm at most $U$
in every class, and
$h\le U^{4/3}\zeta_L(4/3)\le4^{12}U^{4/3}\le H$.
The Euler-factor comparison uses at most twelve primes above
each rational prime, and
$\zeta(4/3)\le1+\int_1^\infty x^{-4/3}dx=4$.
The exact bounds for the three fields are
\[
\begin{array}{c|r|r}
d&U&H\\\hline
-2959&4003722257532096&10666087753422101304760598528\\
538&2737999226454874&6426349247033217033686745088\\
-5918&2049905795856432870&43688295438016918930848449822720.
\end{array}
\tag{G137.22}
\]

The analytic class-number formula has exactly the normalization
\[
 h_{L_d}R_{L_d}
   =\frac{2\sqrt{D(L_d)}}{2^{r_1}(2\pi)^{r_2}}
                \operatorname*{res}_{s=1}\zeta_{L_d}(s).
 \tag{G137.23}
\]
The factors $2^{r_1}$ and $(2\pi)^{r_2}$ therefore differ in
the real-quartic and imaginary-quartic cases as required.
Combining the exact factors above with (G137.8) gives outward
intervals for $r_d=h_{L_d}/m_d$ and proves $r_d>1$.
All three completed runs give the following strict comparisons:
\[
\begin{array}{c|r|c|c}
d&c_d&|r_d-c_d|\text{ less than}&1/H\text{ greater than}\\\hline
-2959&168&8\cdot10^{-49}&9.37\cdot10^{-29}\\
538&22848&2\cdot10^{-47}&1.55\cdot10^{-28}\\
-5918&2688&6\cdot10^{-48}&2.28\cdot10^{-32}.
\end{array}
\tag{G137.24}
\]
The script \texttt{certify\_shared\_quartic.py}, SHA256
\texttt{a5b9d71423f623419d8aa72f860e074f}\allowbreak
\texttt{3fa2e0d807950c78cfc26b211d432ac5},
makes these comparisons on ball objects and preserves every
input hash, phase, theta value, tail, log row and runtime precision
in the three \texttt{analytic\_interval.json} files.

Since $r_d>1$, one has $m_d<h_{L_d}\le H$.
If $h_{L_d}-c_dm_d$ were a nonzero integer, it would force
$|r_d-c_d|\ge1/m_d>1/H$, contradicting (G137.24).
Hence $h_{L_d}=c_dm_d$.  The integers $168,22848,2688$ and
the unit indices in (G137.7) are all prime to $269$.
This proves (G137.2), without asserting any index equals one.

Finally, on the $269$-primary class group of $M_d$, extension
and norm through the reflection-fixed field $L_d$ have
composites $2$ and $1+s$.
The target class group has no $269$-primary part, so $s$ acts
as minus one.  All conjugate reflections do likewise, and
their products are the cubic rotations.  Those rotations act
trivially; the nontrivial cubic class sector is therefore zero.
No trivial-cubic-sector conclusion is used.
\end{proof}

\paragraph{Source and exact certification scope.}
The retained class-number formula and weighted-log convention
are Sutherland, MIT 18.785 Lecture 19 (2019), printed pp.1 and 4,
\url{https://math.mit.edu/classes/18.785/2019fa/LectureNotes19.pdf}.
Primitive Dirichlet special values and functional equations have
the scope of DLMF 25.15.3 and 25.15.5--6,
\url{https://dlmf.nist.gov/25.15}.
The ideal-count isolation and the prime-power Fourier/Poisson
argument are explicitly proved in AC-G135 and re-applied here
with every changed local conductor and signature checked.
The new shared-$269$ different and order-six calculations are
given in Steps 1 and 4, rather than inferred from coprime conductors.
The runtime is python-flint 0.9.0, with its documented outward
ball operations and order-one \texttt{arb.expint}.
No implementation proof, GRH assumption, provisional class-group
completeness, complete unit basis or Lean execution is claimed.

The retained class-field/finite-etale dictionary (SGA1, Expose V,
sections 4 and 7) interprets the result as control of the indicated
cubic sector of unramified abelian covers of
$\operatorname{Spec}\mathcal O_{M_d}$.
A change of residual tuple, local conditions or growing host
still requires its own arithmetic and model comparison.
